% Part: first-order-logic
% Chapter: tableaux
% Section: identity

\documentclass[../../../include/open-logic-section]{subfiles}

\begin{document}

\olfileid{fol}{tab}{ide}

\olsection{\usetoken{P}{tableau} met \usetoken{S}{identity}}

!!^{tableau}s met !!{identity} vereisen aanvullende afleidingsregels.
De regels voor $\eq$ luiden als volgt (waarbij $t$, $t_1$ en $t_2$
gesloten termen zijn):

\begin{defish}
\AxiomC{}
\RightLabel{$\eq$}
\UnaryInfC{\sFmla{\True}{\eq[t][t]}}
\DisplayProof
\hfill
\AxiomC{\sFmla{\True}{\eq[t_1][t_2]}}
\noLine
\UnaryInfC{\sFmla{\True}{!A(t_1)}}
\RightLabel{$\TRule{\True}{\eq}$}
\UnaryInfC{\sFmla{\True}{!A(t_2)}}
\DisplayProof
\hfill
\AxiomC{\sFmla{\True}{\eq[t_1][t_2]}}
\noLine
\UnaryInfC{\sFmla{\False}{!A(t_1)}}
\RightLabel{$\TRule{\False}{\eq}$}
\UnaryInfC{\sFmla{\False}{!A(t_2)}}
\DisplayProof
\end{defish}
Merk op dat $\TRule{\True}{\eq}$ en $\TRule{\False}{\eq}$, in
tegenstelling tot alle andere regels, vereisen dat er al \emph{twee}
!!{signed formula}s op de tak voorkomen, namelijk zowel
$\sFmla{\True}{\eq[t_1][t_2]}$ als $\sFmla{S}{!A(t_1)}$.

\begin{ex}
Als $s$ en $t$ gesloten termen zijn, dan geldt
$\eq[s][t], !A(s) \Proves !A(t)$:
\begin{oltableau}
  [\sFmla{\False}{\formula{A}(t)}, just = \TAss
    [\sFmla{\True}{\eq[s][t]}, just = \TAss
      [\sFmla{\True}{\formula{A}(s)}, just = \TAss
        [\sFmla{\True}{\formula{A}(t)}, just={\TRule{\True}{\eq}[2, 3]}, close]
      ]
    ]
  ]
\end{oltableau}
Dit staat bekend als het substitutieprincipe voor identiteit, ook wel
het principe van Leibniz.

!!^{tableau}s bewijzen dat $\eq$ symmetrisch is, dat wil zeggen dat
$\eq[s_1][s_2] \Proves \eq[s_2][s_1]$:
\begin{oltableau}
  [\sFmla{\False}{\eq[s_2][s_1]}, just = \TAss
    [\sFmla{\True}{\eq[s_1][s_2]}, just = \TAss
      [\sFmla{\True}{\eq[s_1][s_1]}, just = {$\eq$}
        [\sFmla{\True}{\eq[s_2][s_1]}, just = {\TRule{\True}{\eq}[2, 3]}, close]
      ]
    ]
  ]
\end{oltableau}
Hier is regel~$2$ de eerste vereiste
!!{formula}~$\sFmla{\True}{\eq[s_1][s_2]}$ van
$\TRule{\True}{\eq}$. Regel~$3$ is de tweede, van de vorm
$\sFmla{\True}{!A(s_2)}$. Beschouw $!A(x)$ als $\eq[x][s_1]$. Dan is
$!A(s_1)$ gelijk aan $\eq[s_1][s_1]$ en $!A(s_2)$ aan
$\eq[s_2][s_1]$.

Ze bewijzen ook dat $\eq$ transitief is, dat wil zeggen dat
$\eq[s_1][s_2], \eq[s_2][s_3] \Proves \eq[s_1][s_3]$:
\begin{oltableau}
  [\sFmla{\False}{\eq[s_1][s_3]}, just = \TAss
    [\sFmla{\True}{\eq[s_1][s_2]}, just = \TAss
      [\sFmla{\True}{\eq[s_2][s_3]}, just = \TAss
        [\sFmla{\True}{\eq[s_1][s_3]}, just = {\TRule{\True}{\eq}[3, 2]}, close]
      ]
    ]
  ]
\end{oltableau}
In dit !!{tableau} is de eerste vereiste !!{formula} van
$\TRule{\True}{\eq}$ regel~$3$,
$\sFmla{\True}{\eq[s_2][s_3]}$; $s_2$ speelt de rol van~$t_1$ en
$s_3$ die van~$t_2$. De tweede vereiste formule, van de vorm
$\sFmla{\True}{!A(s_2)}$, is regel~$2$. Beschouw $!A(x)$ hier als
$\eq[s_1][x]$. Daardoor wordt $!A(s_2)$ gelijk aan
$\eq[t_1][t_2]$ (d.w.z. regel~$2$) en $!A(s_3)$ aan de
!!{formula}~$\eq[s_1][s_3]$ in de conclusie.
\end{ex}

\begin{prob}
Geef gesloten !!{tableau}s voor de volgende gevallen:
\begin{enumerate}
\item $\sFmla{\False}{\lforall[x][\lforall[y][((x = y \land !A(x))
      \lif !A(y))]]}$
\item $\sFmla{\False}{\lexists[x][(!A(x) \land
    \lforall[y][(!A(y) \lif y = x)])]}$,\\
  $\sFmla{\True}{\lexists[x][!A(x)] \land
    \lforall[y][\lforall[z][((!A(y) \land !A(z)) \lif y = z)]]}$
\end{enumerate}
\end{prob}

\end{document}
